\begin{minipage}{0.58\linewidth}
    Un moteur fait tourner un disque homogène de diamètre ${d=\qty{20}{\cm}}$
    autour d'un axe fixe ($\Delta$) passant par son centre.
    
    On donne la représentation de la variation de l'abscisse angulaire en
    fonction du temps.

    \begin{enumerate}
    \item Quelle est la nature du mouvement de rotation du disque? Justifier
          votre réponse. \textbf{(1pt)}
    \item
          \begin{enumerate}
              \item Déterminer graphiquement la vitesse angulaire $\omega$ et la
                    valeur de l'abscisse angulaire $\theta_{0}$ à $t=0$.
                    \textbf{(1pt)}
              \item Écrire l'équation horaire $\theta(t)$ du mouvement du
                    disque. \textbf{(0,5pt)}
              %\item Déterminer la valeur de la fréquence $f$ du mouvement de
              %      rotation du disque en ($\unit{\Hz}$) puis en
              %      ($\unit{\tr\per\min}$). \textbf{(1pt)}
              %\item Déterminer la valeur de la période $T$ de rotation du
              %      disque. \textbf{(0,5pt)}
          \end{enumerate}
    \end{enumerate}
\end{minipage}
\hfill
\begin{minipage}{0.40\linewidth}
        \flushright
        \begin{tikzpicture}
            \draw[lightgray] (0,0) grid (5.5,5.5);
            \draw[-Latex] (0,0) -- (5.7,0) node[above] {$t~(\unit{\second})$};
            \draw[-Latex] (0,0) -- (0,5.7)
                node[right] {$\theta~(\unit{\radian})$};
            \draw[very thick] (0,1) -- ++(45:6);
            \foreach \x/\l in {0/0,1/0.5,2/1,3/1.5,4/2,5/2.5}
            \draw (\x,3pt) -- (\x,0) node[below] {$\num{\l}$};
            \foreach \y/\l in {0/0,1/12.7,2/25.4,3/38.1,4/50.8,5/63.5}
            \draw (3pt,\y) -- (0,\y) node[left] {$\num{\l}$};
        \end{tikzpicture}
\end{minipage}

\begin{enumerate}
    \setcounter{enumi}{2}
    \item[]
          \begin{enumerate}
              \setcounter{enumii}{2}
              %\item Déterminer graphiquement la vitesse angulaire $\omega$ et la
              %      valeur de l'abscisse angulaire $\theta_{0}$ à $t=0$.
              %      \textbf{(1pt)}
              %\item Écrire l'équation horaire $\theta(t)$ du mouvement du
              %      disque. \textbf{(0,5pt)}
              \item Déterminer la valeur de la fréquence $f$ du mouvement de
                    rotation du disque en ($\unit{\Hz}$) puis en
                    ($\unit{\tr\per\min}$). \textbf{(1pt)}
              \item Déterminer la valeur de la période $T$ de rotation du
                    disque. \textbf{(0,5pt)}
          \end{enumerate}
    \item Donner l'équation horaire de l'abscisse curviligne $s(t)$ d'un point
          du périmètre du disque. \textbf{(1pt)}
    \item
          \begin{enumerate}
              \item Calculer la valeur de $\theta$ à l'instant~:
                    $t=\qty{0,25}{\second}$. \textbf{(0,5pt)}
              \item Quel est le nombre de tours $n$ effectués par le disque à
                    l'instant~: $t=\qty{0,25}{\second}$~? \textbf{(1pt)}
              \item Sachant que un point $M$ du disque a pour vitesse
                    $v=\qty{1,27}{\meter\per\s}$, déterminer la distance qui le
                    sépare de l'axe de rotation. \textbf{(0,5pt)}
          \end{enumerate}
\end{enumerate}
